\documentclass[a4paper,12pt]{article}
\usepackage[T1]{fontenc}
\usepackage{amsmath, graphicx, color, amsthm, tikz}
\usepackage[font=small,labelfont=bf,width=12.5cm]{caption}
\usepackage{algorithmic, algorithm}
%\usepackage[retainorgcmds]{IEEEtrantools}
%\usetikzlibrary{arrows}

\newcommand{\midarrow}{\tikz \draw[-triangle 90] (0,0) -- +(.1,0);}

\addtolength{\textwidth}{30mm}
\addtolength{\hoffset}{-15mm}
\addtolength{\textheight}{30mm}
\addtolength{\voffset}{-15mm}

\newtheorem{theorem}{Theorem}[section]
%\newtheorem{claim}{Claim}[section]
\newtheorem{lemma}[theorem]{Lemma}
\newtheorem{corollary}[theorem]{Corollary}
\newtheorem{proposition}[theorem]{Proposition}
\newtheorem{conjecture}[theorem]{Conjecture}
\newtheorem{question}[theorem]{Question}
\newtheorem{problem}[theorem]{Problem}
\newtheorem{definition}[theorem]{Definition}
\numberwithin{equation}{section}

\newcommand{\inst}[1]{$^{#1}$}
\newcommand{\set}[1]{\ensuremath{\left\{#1 \right\}}}
\newcommand{\red}[1]{\ensuremath{\mathrm{red}(#1)}}
\newcommand{\chis}[1]{\chi_{\textrm{st}}'(#1)}
\newcommand{\comment}[1]{\ $[\![${\scriptsize #1}$]\!]$ \ }

\title{Odd edge-colorings of subdivisions of odd graphs}


\author
{
	Mirko Petru\v{s}evski \thanks{Department of Mathematics and Informatics, Faculty of Mechanical Engineering - Skopje, Republic of North Macedonia. E-Mail: \texttt{mirko.petrushevski@gmail.com}},
	\quad
	Riste \v{S}krekovski\thanks{FMF, University of Ljubljana \& Faculty of Information Studies, Novo mesto, Slovenia. E-Mail: \texttt{skrekovski@gmail.com}}
}


\begin{document}
\maketitle

\begin{abstract}
An odd graph is a finite graph all of whose vertices have odd degrees. Given graph $G$ is decomposable into $k$ odd subgraphs if its edge set can be partitioned into $k$ subsets each of which induces an odd subgraph of $G$. The minimum value of $k$ for which such a decomposition of $G$ exists is the odd chromatic index, $\chi_{o}'(G)$, introduced by Pyber (1991). For every $k\geq\chi_{o}'(G)$, the graph $G$ is said to be odd $k$-edge-colorable. Apart from two particular exceptions, which are respectively odd $5$- and odd $6$-edge-colorable, the rest of connected loopless graphs are odd $4$-edge-colorable, and moreover one of the color classes can be reduced to size $\leq2$. In addition, it has been conjectured that an odd $4$-edge-coloring with a color class of size at most $1$ is always achievable. Atanasov et al. (2016) characterized the class of subcubic graphs in terms of the value $\chi_{o}'(G)\leq4$. In this paper, we extend their result to a characterization of all subdivisions of odd graphs in terms of the value of the odd chromatic index. This larger class $\mathcal{S}$ is of a particular interest as it collects all `least instances' of non-odd graphs.  As a prelude to our main result, we show that every connected graph $G\in \mathcal{S}$ requiring the maximum number of four colors, becomes odd $3$-edge-colorable after removing a certain edge. Thus, we provide support for the mentioned conjecture by proving it for all subdivisions of odd graphs. The paper concludes with few problems for possible further work.
\end{abstract}

\medskip

\noindent \textbf{Keywords:} odd graph, odd edge-coloring, odd chromatic index, subdivision.

%-----------------------------------------------------------------------------------------------------------------------
%-----------------------------------------------------------------------------------------------------------------------

\section{Introduction}

\subsection{Basic terminology}

All considered graphs $G=(V(G),E(G))$ are undirected and finite, loops and parallel edges are allowed. We follow~\cite{BonMur08} for the terminology and notation not defined here. The parameters $n(G)=|V(G)|$ and $m(G)=|E(G)|$ are called \textit{order} and \textit{size} of $G$, respectively. A graph of order $1$ is \textit{trivial}, and a graph of size $0$ is \textit{empty}. A path or cycle is said to be \textit{odd} (resp. \textit{even}) if its size is odd (resp. even.) A path (resp. edge) with endpoints $x$ and $y$ is referred to as an \textit{$x$-$y$ path} (resp. \textit{$x$-$y$ edge}). Given a path $P$ and $x,y\in V(P)$, the $x$-$y$ subpath of $P$ is denoted $xPy$.
For every vertex $v\in V(G)$, $E_G(v)$ denotes the set of edges incident with $v$, and the size of $E_G(v)$ (every loop being counted twice) is the \textit{degree} of $v$ in $G$, with notation $d_{G}(v)$. The maximum (resp. minimum) vertex degree in $G$ is denoted by $\Delta(G)$ (resp. $\delta(G)$).  A graph $G$ is \textit{subcubic} if $\Delta(G)\leq 3$. Each vertex $v$ of even (resp. odd) degree $d_{G}(v)$ is an \textit{even} (resp. \textit{odd}) vertex. In particular, if $d_G(v)$ equals $0$ (resp. $1$), we say that $v$ is an \textit{isolated} (resp. \textit{pendant}) vertex of $G$. Any vertex of degree $d$ is called a \textit{$d$-vertex}. A graph is \textit{even} (resp. \textit{odd})
whenever all its vertices are even (resp. odd).
The set of neighboring vertices of $v\in V(G)$ is denoted by $N_G(v)$. For every $u\in N_G(v)$, the edge set $E_G(u)\cap E_G(v)$ is the \textit{$u$-$v$ bouquet} in $G$, with notation $\mathcal{B}_{uv}$.  The maximum size of a bouquet in $G$ is its \textit{multiplicity}. A graph $G$ is \textit{simple} if it is loopless and of multiplicity at most $1$.

To \textit{split} a vertex $v$ is to replace it by two (not necessarily adjacent) vertices $v'$ and $v''$, and to replace each edge incident to $v$ by an edge incident to either $v'$ or $v''$ (but not both, unless the edge is a loop at $v$), the other end of the edge remaining unchanged. A vertex of positive degree can be split in several ways, so the resulting graph is not unique in general. Another local operation on graph $G$ is to \textit{suppress} a vertex $v$ of degree $d_G(v)=2$. The modified graph $G\%v$ is obtained from $G-v$ by adding an edge between the neighbors of $v$ (the new edge is a link unless $N_G(v)$ is a singleton).
For $X\subseteq V(G)\cup E(G)$, $G-X$ is the subgraph of $G$ obtained by removing $X$; we abbreviate $G-\{x\}$ to $G-x$. Similarly, given a subgraph $H\subseteq G$, $H+X$ is the subgraph of $G$ obtained by adding to $H$ all the vertices and edges from $X$.

 A $v\in V(G)$ is a \textit{cut-vertex} of $G$ if $G-v$ has more (connected) components than $G$. If $V_1,\ldots,V_k$ are vertex sets of the components of $G-v$, then for $i=1,\ldots,k$, the induced subgraph $G[V_i\cup \{v\}]$ is called a \textit{$v$-lobe} of $G$.
 A \textit{block graph} is a connected graph without cut-vertices. Given a graph $G$, any maximal block subgraph is a \textit{block} of $G$. For a block $B$, each vertex $v\in V(B)$ which is not a cut-vertex of $G$ is an \textit{internal vertex} of $B$ (and of $G$). The collection of internal vertices of $B$ is denoted by $\mathrm{Int}_G(B)$.
 If $V(B)$ contains at most one cut-vertex of $G$ then $B$ is an \textit{end-block}. Every vertex $v$ of a block graph $G$ has a neighbor among the internal vertices of each end-block of $G-v$.
 For a connected graph $G$, an edge $e\in E(G)$ is a \textit{bridge} of $G$ if $G-e$ is disconnected. If $vw$ is a bridge and $w$ is not the only neighbor of $v$, then $v$ is a cut-vertex of $G$. The \textit{connectivity} $\kappa(G)$ (resp. \textit{edge-connectivity} $\kappa'(G)$) of a graph $G$ is the minimum size of a subset $S\subseteq V(G)$ (resp. $S\subseteq E(G)$) such that $G-S$ is disconnected or of order $1$.




%----------------------------------------------------------------------------------------------------------------------

\subsection{Odd edge-colorings and odd chromatic index}

An assignment $\varphi: E(G)\to S$ is an \textit{edge-coloring of $G$} with \textit{color set} $S$. We say that $\varphi$ is a $k$-edge-coloring if $|S|\leq k$. Nature of the colors being irrelevant, it is conventional  to use $S=[k]:=\{1,2,\ldots,k\}$ if the color set is of size $k$. For each color $c\in S$, $E_c(G,\varphi)$ denotes the \textit{color class of $c$}, being the set of edges colored by $c$ (i.e., $E_c(G,\varphi)=\varphi^{-1}(c)$); whenever $G$ and $\varphi$ are clear from the context, we denote the color class of $c$ simply by $E_c$.  Given an edge-coloring $\varphi$ and a vertex $v$ of $G$, we say the color $c$ \textit{appears at $v$} if $E_c\cap E_G(v)\neq \emptyset$. Any decomposition $\{H_1,\ldots, H_k\}$ of $G$ can alterably be interpreted as an $k$-edge-coloring of $G$ for which the color classes are $E(H_1),\ldots, E(H_k)$.

An \textit{odd edge-coloring} of a given graph $G$ is an edge-coloring such that each nonempty color class induces an odd subgraph. In other words, at each vertex $v$, for any appearing color $c$ the degree $d_{G[E_c]}(v)$ is odd. Equivalently, an odd edge-coloring can be seen as a decomposition of $G$ into
(edge-disjoint) odd subgraphs. As usual, we are most interested in the least number of colors necessary to create such a coloring.
An odd edge-coloring of $G$ using at most $k$ colors is referred to as an \textit{odd $k$-edge-coloring}, and then we say that $G$ is
\textit{odd $k$-edge-colorable}. Whenever $G$ is odd edge-colorable, the \textit{odd chromatic index} $\chi_{o}'(G)$ is
defined to be the minimum integer $k$ for which $G$ is odd $k$-edge-colorable.

It is obvious that a necessary and sufficient condition for odd edge-colorability of $G$ is the absence of vertices incident only to loops. Apart from this, the presence of loops
does not influence the existence nor changes the value of the index $\chi_{o}'(G)$. Therefore, the class of loopless graphs comprise a natural framework for the study of the odd chromatic index.

\begin{figure}[ht!]
	$$
		\includegraphics{wheel}
	$$
	\caption{A simple graph with odd chromatic index equal to $4$.}
	\label{fig:pyb}
\end{figure}

As a notion, odd edge-coloring was introduced by Pyber in his survey on graph coverings~\cite{Pyb91}.
The mentioned work considers simple graphs and (among other results) contains a proof of the following.

\begin{theorem}[Pyber, 1991]
	\label{thm:pyb}
	For every simple graph $G$, it holds that
	$\chi_o'(G) \le 4$.
\end{theorem}


Pyber remarked that the established upper bound is realized by the wheel on four spokes $W_4$ (see Fig.~\ref{fig:pyb}). However, this upper bound of four colors does not apply to the class of all looplees graphs $G$. For instance, Fig.~\ref{fig:shan} depicts four graphs with the following characteristic property: each of their odd subgraphs is of order $2$ and size $1$, i.e., a copy of $K_2$. Consequently, for each of these graphs the odd chromatic index equals the size.

\begin{figure}[ht!]
	$$
		\includegraphics{shannons}
	$$
	\caption{Four Shannon triangles (the smallest one of each type).}
	\label{fig:shan}
\end{figure}

As defined in~\cite{LuzPetSkr15}, a \textit{Shannon triangle} is a loopless graph on three pairwise adjacent vertices. And if $p,q,r$ are parities of the sizes of its bouquets in non-increasing order, with $2$ (resp. $1$) denoting an even-sized (resp. odd-sized) bouquet, then $G$ is a Shannon triangle of \textit{type $(p,q,r)$}. For example, Figure~\ref{fig:shan} depicts (from left to right) the smallest in terms of size Shannon triangle of type $(1,1,1)$, $(2,1,1)$, $(2,2,1)$, and $(2,2,2)$, respectively. It is straightforward that if $G$ is a Shannon triangle of type $(p,q,r)$, then
\begin{equation}
    \label{eqn:shan}
\chi'_{o}(G)=p+q+r\,.
\end{equation}
The main result of~\cite{LuzPetSkr15} is that six colors suffice for an odd edge-coloring of any connected loopless graph. Furthermore, it characterizes when six colors are necessary.
\begin{theorem}
    \label{char:6}
    For every connected loopless graph $G$, it holds that
	$\chi_{o}'(G) \leq 6$.
	Moreover, equality is attained if and only if $G$ is a Shannon triangle of type $(2,2,2)$.
\end{theorem}

Recently, the following improvement of Theorems~\ref{thm:pyb} and~\ref{char:6} has been shown in~\cite{Pet18}:

\begin{theorem}
 \label{odd 4-edge-colorability}
 Let $G$ be a connected loopless graph that is not a Shannon triangle of type $(2,2,1)$ or $(2,2,2)$. Then $G$ admits an odd edge-coloring with color set $\{1,2,3,4\}$ such that the color class $E_4$ satisfies two additional conditions:
\begin{enumerate}
\item[$(i)$] $|E_4|\in\{0,1,2\}$;
\item[$(ii)$] if $\mathcal{B}_{xy}\cap E_4\neq \emptyset$ then another common color (besides $4$) appears at $x$ and $y$;
\end{enumerate}
and in case of $|E_4|=2$, the pair of edges colored by $4$ are at distance $2$ (i.e., second-neighbors in the line graph).
\end{theorem}

It is not known whether there exists a connected graph $G$ with $\chi'_o(G)=4$ that does not admit an odd $4$-edge-coloring with a color class of size $1$. In this regard, the following has been conjectured in~\cite{PetSkr21}.

\begin{conjecture}
    \label{conj:finiteness}
Every connected graph $G$ with $\chi'_o(G)=4$ becomes odd $3$-edge-colorable by removing a particular edge.
\end{conjecture}

Concerning odd $2$-edge-colorability of graphs, Kano et al.~\cite{KanKatVar18} have shown the following.

\begin{theorem}
    \label{Kano}
The decision problem whether a given graph $G$ is odd $2$-edge-colorable is solvable in polynomial time. Moreover, in the affirmative case, such a coloring can be found in polynomial time.
\end{theorem}

In view of Theorem~\ref{odd 4-edge-colorability}, the decision problem whether a given graph $G$ is odd $4$-edge-colorable is also solvable in polynomial (in fact, linear) time. Moreover, its proof can be used as an efficient algorithm for exhibiting such a coloring. The analogous complexity questions regarding odd $3$-edge-colorability of general graphs are still open. Nevertheless, these questions have been answered (in the affirmative) for subcubic graphs. Namely, a complete characterization of the class of loopless subcubic graphs in terms of their odd chromatic index was obtained in~\cite{AtaPetSkr16} through the following:

\begin{theorem}
    \label{thm:characterizationsubcubic}
Let $G$ be a connected loopless subcubic graph. Then
\begin{equation*}
\chi'_o(G)=
\begin{cases}
0 & \text{\quad if\, } G \text{ is empty}\,;\\
1 & \text{\quad if \,} G \text{ is odd} \,;\\
2 & \text{\quad if \,} G \text{ has 2-vertices, with an even number of them on each cycle}\,;\\
4 & \text{\quad if \,} G \text{ is obtainable from a cubic bipartite graph by a single edge subdivision}\,;\\
3 & \text{\quad otherwise}\,.
\end{cases}
\end{equation*}
\end{theorem}

In this paper we focus on subdivisions of odd graphs, which are in some sense the least non-odd graphs among all. Let us denote this collection by $\mathcal{S}$, and similarly, let $\mathcal{O}$ be the class of odd graphs, with the understanding that $\mathcal{O}\subset\mathcal{S}$. Our main result, see Theorem~\ref{thm} in Section~$3$, is a characterization of the members of $\mathcal{S}$ in terms of the value of their odd chromatic index. Thus, we achieve a generalization of Theorem~\ref{thm:characterizationsubcubic}, and at the same time answer a question raised by the end of \cite{Pet18}. Our findings here also provide support for  Conjecture~\ref{conj:finiteness} over the class $\mathcal{S}$.

\medskip

The rest of the article is divided into three sections. In the next one, as a preliminary, we collect several `easy' results (most of them previously known). Section~$3$ is devoted to a derivation of our main result - a characterization of $\mathcal{S}$ in terms of the value of the odd chromatic index. The final section briefly conveys some possible directions for further related study.

%----------------------------------------------------------------------------------------------------------------------
%----------------------------------------------------------------------------------------------------------------------

\section{Preliminaries}

If $H\subseteq G$ and the underlying (host) graph $G$ is clear from the context, then the edge-complement of the subgraph $H$, denoted  $\hat{H}$, is the spanning subgraph $\hat{H}=G-E(H)$.
A \textit{co-forest} in $G$ is a subgraph whose edge-complement is a forest.
For a graph $G$, let $T$ be an even-sized subset of $V(G)$. Following~\cite{BonMur08}, a spanning subgraph $H$ of $G$ is said to be a \textit{$T$-join} of $G$ if $d_{H}(v)$ is odd for all $v\in T$ and even for all $v\in V(G)\setminus T$. For instance, if $P\subseteq G$ is a path with endpoints $x$ and $y$, the spanning subgraph of $G$ with edge set $E(P)$ is an $\{x,y\}$-join of $G$. As another example, every even spanning subgraph is an $\emptyset$-join of $G$. Observe that the symmetric difference of an $S$-join and a $T$-join is an $S\oplus T$-join. (We shall use $\oplus$ to denote both the symmetric difference operation on spanning subgraphs and on sets.) Hence, the symmetric difference $H\oplus K$ of a $T$-join $H$ and a spanning even subgraph $K$ of $G$, is a $T$-join. In particular, removal (resp. addition) of the edges of an edge-disjoint cycle from (resp. to) a $T$-join, produces another $T$-join. Therefore, if a $T$-join of $G$ exists, there also exists such a forest (resp. co-forest). By the handshake lemma, necessary for the existence of a $T$-join is that the intersection of $T$ with the vertex set of every component of $G$ is even-sized, and a straightforward implementation of the  above mentioned facts (see~\cite{Sch03}) is that this condition also suffices. Consequently, given a connected graph $G$ and an even-sized subset $T$ of $V(G)$,
\begin{enumerate}
\item[$(1)$] there exists a $T$-join of $G$ that is a forest;
\item[$(2)$] there exists a $T$-join of $G$ that is a co-forest;
\item[$(3)$] additionally, if $G$ is of even order, then it contains a spanning odd co-forest.
\end{enumerate}


\smallskip

 An edge-coloring $\varphi$ is said to be \textit{odd} (resp. \textit{even}) \textit{at} vertex $v$ if each color appearing  at $v$ is odd (resp. even). Similarly, we say that $\varphi$ is odd (resp. even) \textit{away from $v$} if $\varphi$ is odd (resp. even) at every vertex $w\in V(G)\backslash\{v\}$, with no assumptions about the behavior of $\varphi$ at $v$ being made.
The following useful result appears in~\cite{Mat06,Pet18}.

\begin{proposition}
    \label{forest}
    Let $v$ be a vertex in a forest $F$. Any local coloring of $E_F(v)$ which uses at most two colors extends to a $2$-edge-coloring of $F$ that is odd away from $v$. In particular, $F$ is odd $2$-edge-colorable.
\end{proposition}


An immediate consequence of Proposition~\ref{forest} is the result below, which concerns a graph all of whose cycles (if any) share a vertex.
 \begin{proposition}
    \label{consequences:forest}
    If $v$ is a vertex of a graph $G$ such that $G-v$ is a forest, then $G$ admits a $2$-edge-coloring that is odd away from $v$. Additionally, if $d_G(v)$ is odd, then $G$ admits an edge-coloring with color set $\{1,2\}$ that is odd away from $v$ and the color $1$ (resp. $2$) is odd (resp. even) at $v$.
 \end{proposition}

 \begin{proof}
 We may assume that $G$ is loopless. It suffices to prove the first part. Split $v$ into $k=d_G(v)$ pendant vertices $v_1,\ldots,v_k$ in order to obtain a forest $F$. By Proposition~\ref{forest}, $F$ admits an odd $2$-edge-coloring. Re-identify $v_1,\ldots,v_k$ into $v$ while keeping the colors on all edges. We thus regain $G$ and a required edge-coloring.
 \end{proof}


As observed in~\cite{Pyb91}, the odd $2$-edge-colorability of forests implies odd $3$-edge-colorability of connected graphs of even order, which in turn yields odd $3$-edge-colorability of graphs with edge-connectivity $1$.
\begin{proposition}
     \label{even order}
If $G$ is a connected graph such that $n(G)$ is even or $\kappa'(G)=1$ then $\chi_{o}'(G)\leq 3$.
\end{proposition}



The next result can be used to characterize odd $2$-edge-colorability of unicyclic graphs.

\begin{proposition}
Let $G$ be a unicyclic graph, and let $C\subseteq G$ be the (unique) cycle. Then $\chi'_o(G)\leq3$. Moreover, the upper bound is attained if and only if the following two conditions hold simultaneously:
\begin{itemize}
\item[$(i)$] $\{v\in V(C):d_G(v)=2\}$ is odd-sized;
\item[$(ii)$] $\{v\in V(C):0\equiv_2d_G(v)\neq2 \}$ is empty.
\end{itemize}
\begin{proof}
If $G=C$ then condition $(ii)$ is clearly met, and the statement is trivially true (as by then the notions `odd edge-coloring'  and `proper edge-coloring' become equivalent). Assuming $G\neq C$, let $S=\{v\in V(C):d_G(v)\neq2\}$ and $\hat{S}=V(C)\backslash S$. Denote by $S'$ and $S''$, respectively, the subsets of $S$ comprised of those vertices $v$ for which the degree $d_G(v)$ is odd or even. Observe that for every $v\in S'$, a coloring of $E_C(v)$ extends to an odd $2$-edge-coloring of $E_G(v)$ if and only if it is monochromatic. Contrarily, for every $v\in S''$ each coloring of $E_C(v)$ extends to an odd $2$-edge-coloring of $E_G(v)$. Consequently, in view of Proposition~\ref{forest}, a given $k$-edge-coloring of $E(C)$ extends to an odd $k$-edge-coloring of $G$ if and only if it is dichromatic at each $v\in \hat{S}$ and monochromatic at each $v\in S'$. In particular, it is readily implied that $\chi'_o(G)\leq3$.

Additionally, if condition $(ii)$ fails to hold, then a color can be saved; indeed, simply select a vertex $v\in S''$, take a $2$-edge-coloring of $C$ which is monochromatic at each vertex from $S'$ and dichromatic at each vertex from  $V(C)\backslash (S'\cup \{v\})$; by the above observation, such a coloring of $E(C)$ extends to an odd $2$-edge-coloring of $G$. On the other hand, assuming $(ii)$, odd $2$-edge-colorability of $G$ is equivalent to the existence of a $2$-edge-coloring of $C$ that is dichromatic precisely at each vertex of $\hat{S}$. The latter is clearly equivalent to $(i)$.
\end{proof}
\end{proposition}
\begin{corollary}
    \label{unicycle}
Let $G$ be a connected unicyclic graph, $C\subseteq G$ be the (unique) cycle and $\{v\in V(C):d_G(v) \text{ is odd}\}=\emptyset$. Then $\chi'_o(G)\leq2$ unless $G=C$ is an odd cycle.
\end{corollary}


 We end the preliminaries with two more already known results (proofs can be found in~\cite{Pet18}).

\begin{proposition}
    \label{co-forest}
In a connected loopless graph $G$, let $v$ be an internal vertex, $e\in E_G(v)$, and $T\subseteq V(G)$ be even-sized. Then there exists a $T$-join $H$ of $G$ which is a co-forest, such that $E_{\hat{H}}(v)\subseteq\{e\}$.
\end{proposition}

If additionally the graph $G$ from Proposition~\ref{co-forest} is of even order, by setting $T=V(G)$, we derive the following.
\begin{corollary}
    \label{odd co-forest}
   In a connected loopless graph $G$ of even order, let $v$ be an internal vertex and $e\in E_G(v)$. Then there exists a spanning odd co-forest $H$ of $G$ such that $E_{\hat{H}}(v)\subseteq\{e\}$.

\end{corollary}

%----------------------------------------------------------------------------------------------------------------------
%----------------------------------------------------------------------------------------------------------------------

\section{Subdivisions of odd graphs}

Recall that $\mathcal{S}$ denotes the class of subdivisions of odd graphs. The following proposition is an overture to our subsequent study of the class $\mathcal{S}$ in terms of the value of the odd chromatic index. The final product of the study, our main result here, shall be formulated by the end of the article. As a warm-up, we commence by showing that four colors always suffice for an odd edge-coloring of any member of $\mathcal{S}$. Moreover, if necessary at all, the fourth color can be reduced to a single appearance.
\begin{proposition}
    \label{4-edge-colorability}
Let $G\in \mathcal{S}$ be connected and non-trivial, let $v$ be a $2$-vertex of $G$ and $e\in E_G(v)$. Then $G$ admits an odd edge-coloring with color set $\{1,2,3,4\}$ such that the color class $E_4\subseteq \{e\}$. Moreover, if $\chi'_o(G)=4$ then it holds that:
 \begin{itemize}
 \item [$(i)$] Every $2$-vertex is internal;
 \item [$(ii)$] No $2$-vertices are adjacent.
 \end{itemize}
\end{proposition}
\begin{proof}
If there happens to exist a cut-vertex $u$ of $G$ such that $d_G(u)=2$, then $u$ must be incident with two bridges. Consequently, Proposition~\ref{even order} yields odd $3$-edge-colorability of $G$. Assuming $(i)$, the vertex $v$ is internal. Consequently, the graph $G-e$ is connected and of minimum degree $1$, so it admits an odd edge-coloring with color set $\{1,2,3\}$. By assigning the color $4$ to the edge $e$ we obtain the promised coloring of $E(G)$. This proves the first part and, in addition, confirms $(i)$ whenever $\chi'_o(G)=4$.

As for $(ii)$, still assuming $\chi'_o(G)=4$, suppose there is an edge $f$ whose endpoints $u$ and $w$ are $2$-vertices. Let $g_u$ and $g_w$ be the other edges (besides $f$) incident with $u$ and $w$, respectively. Since $g_u$ and $g_w$ are pendant edges in the graph $G-f$, Proposition~\ref{even order} guarantees that there is an odd edge-coloring $\varphi$ of $G-f$ with color set $\{1,2,3\}$. Extend $\varphi$ to $E(G)$ by assigning $f$ with a color from $\{1,2,3\}\backslash\{\varphi(g_u),\varphi(g_w)\}$. This completes an odd $3$-edge-coloring of $G$, a contradiction.\end{proof}

Note that the first part of Proposition~\ref{4-edge-colorability} supports Conjecture~\ref{conj:finiteness}. For our intended characterization of the members of the class $\mathcal{S}$ in terms of their odd chromatic index, let us denote by $\mathcal{S}_i$ $(i=1,2,3,4)$ the subclass consisting of those $G\in\mathcal{S}$ having $\chi'_o(G)=i$. Clearly, $\mathcal{S}_1$ comprises the class of odd graphs $\mathcal{O}$. At the other end of the spectrum, the second part of Proposition~\ref{4-edge-colorability} gives a pair of necessary conditions for membership into $\mathcal{S}_4$, or equivalently, describes two sufficient conditions for odd $3$-edge-colorability of a subdivision of an odd graph. The following result provides another such condition (which shall be useful on more than one occasion later on in this section).

\begin{proposition}
    \label{sufficient condition 3-edge-colorability}
Let $G\in \mathcal{S}$ be a connected graph, and $C\subseteq G$ be a cycle passing through a $2$-vertex $v$ of $G$. If the cycle $C$ is even or it passes through another $2$-vertex, then $G$ is odd $3$-edge-colorable.
\end{proposition}
\begin{proof}
 By Propositions~\ref{even order} and~\ref{4-edge-colorability}~$(ii)$, respectively, we may assume that the order $n(G)$ and the degree of each neighbor of $v$ are odd. If the neighborhood $N_G(v)$ happens to be a singleton, in which case  there is a cycle $C$ of length $2$ that goes through $v$, the odd $3$-edge-coloring of $G$ is immediate. Indeed, since $G-v$ is connected and of even order, simply take an odd co-forest of $G-v$ and its edge-complement in regard to $G$: the former graph is odd $1$-edge-colorable, whereas the latter is a unicyclic graph that is odd $2$-edge-colorable (as Corollary~\ref{unicycle} applies).

Thus, assuming $N_G(v)=\{u,w\}$, let $P=C-v$ be a $u$-$w$ path that (by assumption) is even or passes through another $2$-vertex of $G$. For an arbitrary odd co-forest $H$ of $G-v$, consider the forest $\hat{H}=G-v-E(H)$. Note that the vertices of odd degree in $\hat{H}$ are precisely $u,w$ and the $2$-vertices of $G$ that are within $G-v$. In particular, every $2$-vertex $\neq v$ is a pendant vertex of the forest $\hat{H}$. Looking at $G-E(H)=\hat{H}+\{uv,vw\}$, there are two possibilities: this graph is either a forest (if $u,w$ do not share a component of $\hat{H}$), or it is a unicyclic graph such that no vertex of the cycle has an odd degree.  Therefore, in view of Corollary~\ref{unicycle}, the graph $\hat{H}+\{uv,vw\}$ is odd $2$-edge-colorable (and consequently $G$ is odd $3$-edge-colorable), unless it is always the case that a component of $\hat{H}$ is an odd $u$-$w$ path, say $Q$.

Supposing the latter, it is essential to note that $P\neq Q$. Namely, if the cycle $C$ happens to be even, then $P$ and $Q$ are of different parities; and if $C$ passes through a $2$-vertex of $G$ contained within $G-v$, then so does $P$ but not $Q$ (as every such $2$-vertex of $G$ is a pendant vertex of $\hat{H}$). Let $Q_u$ and $Q_w$ be, respectively, the components of $u$ and $w$ in $P\cap Q$. Notice that these two are disjoint paths having as an endpoint $u$ and $w$, respectively. Thus $Q_u\cup Q_w\subset P$. Now consider an odd co-forest $H$ of $G-v$ that maximizes $Q_u\cup Q_w$ in respect of $\subseteq$. Say $e_u$ and $e_w$ are the first and the last edge lying outside $Q_u\cup Q_w$ on a traversing of $P$ from $u$ to $w$.

 \begin{figure}[ht!]
	$$
		\includegraphics[scale=0.8]{PQH}
	$$
	\caption{The path $P$ and the components of $\hat{H}$ (all are shaded, apart from $Q$). The edges of $P$ are depicted as heavier, and those of $Q_u\cup Q_w$ as fat. On a traversal of $P$ from $u$ to $w$, the arrows notify the first embarkment and last disembarkment of $P-E(Q)$. Since this happens precisely along the edges $e_u$ and $e_w$, respectively, these two cannot be cycle edges of $\hat{H}\oplus P\oplus Q$. The same is obviously true for any edge of $Q_u\cup Q_w$.}
	\label{fig:PQH}
\end{figure}

Take the symmetric difference $H\oplus P\oplus Q$, with both $P$ and $Q$ seen as spanning subgraphs of $G-v$. The obtained graph is clearly another odd factor of $G-v$, though not necessarily a co-forest. Nevertheless, the important feature (which is preserved in passing from $H$ to $H\oplus P\oplus Q$) is that no edge from $Q_u\cup\{e_u\}\cup \{e_w\}\cup Q_w$ belongs to a cycle contained entirely in the edge-complement of $H\oplus P\oplus Q$ with reference to $G-v$ (cf. Figure~\ref{fig:PQH}). Therefore, by taking as $H'$ any maximal odd factor of $G-v$ subjected to the condition $H\oplus P\oplus Q\subseteq H'$, we are guaranteed  that $H'$ is a spanning odd co-forest of $G-v$ and it is also edge-disjoint from $Q_u\cup Q_w$. Moreover, if $Q'$ denotes the odd $u$-$w$ path that (according to our supposition) constitutes a component of the forest $\hat{H'}=G-v-E(H')$, and $Q'_u,Q'_w$ are defined analogously as before, then $Q_u\subset Q'_u$ and $Q_w\subset Q'_w$. Indeed, as $Q_u\cup Q_w\subseteq Q'$ and $Q'$ constitutes a path, it must be that $Q_u\cup\{e_u\}\subseteq Q'_u$ and similarly $Q_w\cup\{e_w\}\subseteq Q'_w$. Consequently, $Q_u\cup Q_w\subset Q'_u\cup Q'_w$, contradicting our choice of $H$. This settles the proposition.
\end{proof}



Our next result concerns odd $2$-edge-colorability of subdivisions of odd graphs, and thus yields a structural characterization of $\mathcal{S}_2$.

\begin{proposition} The following statements are equivalent for every graph $G\in \mathcal{S}$:
\begin{itemize}
\item[$(i)$] $\chi'_o(G)\leq2$;
\item[$(ii)$] For every cycle $C$ of $G$ the set $\{v:v\in V(C) \text{ and } d_G(v)=2\}$ is even-sized.
\end{itemize}
\end{proposition}
\begin{proof}
Let $G_0$ be the underlying odd graph, that is, a graph whose subdivision is $G$. Noting that $G_0$ is well-defined, distinguish between \textit{even} and \textit{odd} edges $e\in E(G_0)$ depending on the parity of the number of $2$-vertices of $G$ that have been introduced along $e$. The corresponding subsets of $E(G_0)$ are $E_{\mathrm{even}}$ and $E_{\mathrm{odd}}$. Let $P_e$ be the path (or cycle) in $G$ arising from the edge $e$; thus, $e$ is even if and only if the length $\ell(P_e)$ is odd. Define a graph $H$ whose vertices correspond to the components of $G_0-E_{\mathrm{odd}}$; for each $e\in E_{\mathrm{odd}}$, we put an edge in $H$ whose endpoints are the components of $G_0-E_{\mathrm{odd}}$ containing the endpoints of $e$. In other words, $H$ is the graph obtained from $G_0$ by contracting every component of $G_0-E_{\mathrm{odd}}$. The equivalence of $(i)$ and $(ii)$ shall be established via a third statement:
\begin{itemize}
\item[$(iii)$] $H$ is bipartite.
\end{itemize}
Namely, we proceed to show the following string of implications: $(i)\Rightarrow (ii)\Rightarrow(iii)\Rightarrow(i)$.
\begin{itemize}
\item $(i)\Rightarrow (ii)$: Considering an odd $2$-edge-coloring $\varphi$ of $G$, let $C$ be an arbitrary cycle in $G$. For any $v\in V(C)$, the set $E_C(v)$ is dichromatic under $\varphi$ if and only if $d_G(v)=2$. Therefore, $\{v:v\in V(C) \text{ and } d_G(v)=2\}$ is of even size.
\item $(ii)\Rightarrow(iii)$: From a cycle $C$ in $H$, we can obtain a cycle $C_0$ in $G_0$ as follows. For $v\in V(C)$, let $e,f$ be the edges of $C$ incident to $v$ (not necessarily distinct), and let $x,y$ be the endpoints of $e,f$ in the component of $G_0-E_{\mathrm{odd}}$ corresponding to $v$. We expand $v$ into an $x$-$y$ path in that component. Since $C$ visits each vertex at most once, the resulting $C_0$ is a cycle in $G_0$. Let $C'$ be the cycle in $G$ obtained by replacing each $e\in E(C_0)$ with $P_e$. By construction, we have $|\{v:v\in V(C') \text{ and } d_G(v)=2\}|\equiv_2\ell(C)$, where $\ell(C)$ denotes the length of $C$. Hence, $(ii)$ implies that $\ell(C')$ is even. Consequently, $H$ is bipartite.
\item $(iii)\Rightarrow(i)$: Let $X,Y$ be a bipartition of $H$. To every edge $e\in E_{\mathrm{even}}$ assign a color from the set $\{1,2\}$ as follows. If the component of $G_0-E_{\mathrm{odd}}$ containing $e$ represents a vertex in $X$, then color $e$ by $1$; contrarily, color $e$ by $2$. This readily gives a partial $2$-edge-coloring of $G$: namely, for every $2$-vertex $v$ of $G$ lying on an even edge $e$ of $G_0$, color $E_G(v)$ by the color assigned to $e$. The uncolored part of $E(G)$ is $\bigcup_{e\in E_{\mathrm{odd}}} E(P_e)$. Since $e\in E_{\mathrm{odd}}$ has an endpoint within each partite set, the path $P_e$ admits a proper $2$-edge-coloring such that its edge `incident' with $X$ (resp. $Y$) receives color $1$ (resp. $2$). This completes an odd $2$-edge-coloring of $G$.
\end{itemize}
\end{proof}

The above proof provides more than just another glimpse at the beautiful idea of proving an equivalence of two given assertions using a third (not obvious) one, which is equivalent to each of them. Namely, as the fulfillment of statement $(iii)$ can be efficiently checked, our characterization of odd $2$-edge-colorability within $\mathcal{S}$ is good and, in the affirmative, such a coloring can be found in polynomial time.

\begin{corollary}
    \label{odd 2-edge-colorability}
Let $G\in\mathcal{S}$. Then $\chi'_o(G)=2$ if and only if $G\notin\mathcal{O}$ and for every cycle $C$ of $G$ the set $\{v:v\in V(C) \text{ and } d_G(v)=2\}$ is even-sized.
\end{corollary}

In the remainder of the paper we provide a structural characterization of the class $\mathcal{S}_4$. The next result shall allow us to confine to $2$-connected graphs.

 \begin{proposition}
    \label{block tree}
 If $G\in\mathcal{S}$ is a connected graph, then  $\chi'_o(G)=4$ if and only if every block of $G$ belongs to $\mathcal{S}_4$ and for every cut-vertex $v$ there is a unique block $B$ such that $d_B(v)$ is odd.

\end{proposition}
\begin{proof}
The essential part of our proof is to establish property $(P)$ below, which sheds some light on the structure of graphs $G\in\mathcal{S}$ of connectivity $\kappa(G)=1$ that require four colors for an odd edge-coloring.

\medskip

\noindent $(P)$ \textit{Let $G\in\mathcal{S}$ be connected and $v$  a cut-vertex of $G$. If $G_1,\ldots,G_k$ are the $v$-lobes of $G$, then the following statements are equivalent:}
\begin{itemize}
\item[$(i)$]  \textit{$\{G_1,G_2,\ldots,G_k\}\subseteq\mathcal{S}_4$  and there is a unique $j$ such that $d_{G_j}(v)$ is odd;}
\item[$(ii)$] \textit{$G\in\mathcal{S}_4$.}
\end{itemize}


Notice that, once the equivalence stated in $(P)$ is verified, the proposition may be derived by inducting on the number $t$ of cut-vertices in $G$. Namely, the case $t=0$ is trivial, whereas the case $t=1$ follows immediately from $(P)$. For $t>1$, consider an end-block of $G$, say $B_0$, and let $v_0$ be the corresponding cut-vertex of $G$. By applying $(P)$ to the pair $G,v_0$ we deduce that $G\in\mathcal{S}_4$ if and only if both $B_0,G-\mathrm{Int}(B_0)\in\mathcal{S}_4$. Now the inductive hypothesis applies to $G-\mathrm{Int}(B_0)$.

\medskip

In what follows, we verify $(P)$ by proving both $(i)\Rightarrow(ii)$ and $(ii)\Rightarrow(i)$.
First we show the easy part, which is direction $(i)\Rightarrow(ii)$. Arguing by contradiction, in view of Proposition~\ref{4-edge-colorability}, assume $(i)$ and consider an odd $3$-edge-coloring $\varphi$ of $G$. For every $i=1,\ldots,k$, let $\varphi_i$ be the restriction of $\varphi$ to $E(G_i)$, i.e., $\varphi_i=\varphi_{|E(G_i)}$. Since $\chi'_o(G_i)=4$, the edge-coloring $\varphi_i$  is even at $v$. Consequently, whenever $d_{G_i}(v)=2$ the edge set $E_{G_i}(v)$ is monochromatic under $\varphi_i$. However, then $\varphi_j$ must be an odd $3$-edge-coloring of $G_j$, a contradiction.

\medskip

Now let us prove direction $(ii)\Rightarrow(i)$. Assuming $(ii)$, note that degree $d_G(v)$ is odd by Proposition~\ref{4-edge-colorability}. We break the argument into several claims  which eventually lead to $(i)$.

\bigskip

\noindent \textbf{Claim 1.} \textit{If $G'\cup G''=G$ and $G'\cap G''=\{v\}$, then orders $n(G'),n(G'')$ are odd.}

\medskip

\noindent By Proposition~\ref{even order}, $n(G)$ is odd. Suppose $n(G'),n(G'')$ are even. Then there exist spanning odd co-forests $H'$ and $H''$ in $G'$ and $G''$, respectively. Denote $F=G-E(H'\cup H'')$. Note that $F$ is a forest and $d_F(v)$ is odd. This enables construction of an odd $3$-edge-coloring of $G$ as follows: color $E(H')$ by $1$ and $E(H'')$ by $2$; color $E_F(v)$ by $3$; extend the coloring of $E_{F\cap G'}(v)$ to an edge-coloring of the forest $F\cap G'$ with color set $\{2,3\}$ that is odd away from $v$; similarly, extend the coloring of $E_{F\cap G''}(v)$ to an edge-coloring of the forest $F\cap G''$ with color set $\{1,3\}$ that is odd away from $v$. The obtained contradiction proves our point.\hfill$\diamond$

\bigskip

Consequently, every $v$-lobe of $G$ has an odd order. We show next that none of the $v$-lobes is odd $3$-edge-colorable.

\bigskip

\noindent \textbf{Claim 2.} \textit{$\chi'_o(G_i)=4$ for every $i=1,2,\ldots,k$.}

\medskip

\noindent Supposing a $v$-lobe $G_i$ is odd $3$-edge-colorable, take an odd edge-coloring $\varphi_i$ of $G_i$ with color set $\{1,2,3\}$ such that the color $1$ appears (resp. does not appear) on $E_{G_i}(v)$ if $d_{G_i}(v)$ is odd (resp. even). Let $\bar{G_i}=\bigcup\{G_j:j=1,2,\ldots,k \text { and } j\neq i\}$. Since $n(G_j)$ is odd for every $j\neq i$, there exists an odd co-forest $K$ in $\bar{G_i}-v$. (Indeed, simply take an odd co-forest $K_j$ in every $G_j-v$ and then form the union of these $K_j$'s.) Color the edge set $E(K)$ by $3$. Use the color $1$ on all of $E_{\bar{G_i}}(v)$ and (by Proposition~\ref{forest}) extend to an edge-coloring of $\bar{G_i}-E(K)$ with color set $\{1,2\}$ that is odd away from $v$. This completes a $3$-edge-coloring $\bar{\varphi_i}$ of $\bar{G_i}$ that is odd away from $v$. However, then $\varphi_i\cup\bar{\varphi_i}$ is an odd $3$-edge-coloring of $G$, a contradiction.
\hfill$\diamond$

\bigskip

Proceed by showing that each $v$-lobe is a subdivision of an odd graph.

\bigskip

\noindent \textbf{Claim 3.} \textit{$G_i\in \mathcal{S}$ for every $i=1,2,\ldots,k$.}

\medskip

\noindent Supposing the opposite, let $G_i$ be a $v$-lobe such that $d_{G_i}(v)$ is even and $d_{G_i}(v)\geq4$. Once again, consider the subgraph $\bar{G_i}=\bigcup\{G_j:j=1,2,\ldots,k \text { and } j\neq i\}$. By Claim~1, every component of $\bar{G_i}-v$ is of even order, and thus there is an odd co-forest $K$ in $\bar{G_i}-v$. Note that the edge-complement $\hat{K}=\bar{G_i}-E(K)$ is such that $\hat{K}-v$ is a forest, and moreover the degree $d_{\hat{K}}(v)=d_{\bar{G_i}}(v)$ is odd. Color $E(K)$ by $3$. By Proposition~\ref{consequences:forest},  $\hat{K}$ admits an edge-coloring with color set $\{1,2\}$ that is odd away from $v$ and under which the color $1$ (resp. $2$) is odd (resp. even) at $v$. This furnishes an edge-coloring $\bar{\varphi_i}$ of $\bar{G_i}$ with color set $\{1,2,3\}$ that is odd away from $v$, whereas at $v$ the color $1$ is odd, the color $2$ is even, and the color $3$ does not appear. We shall prove the existence of an edge-coloring $\varphi_i$ of $G_i$ with color set $\{1,2,3\}$ that is also odd away from $v$, and such that at $v$ both the colors $2$ and $3$ are odd while the color $1$ is even. (Then  $\varphi_i\cup\bar{\varphi_i}$ would be an odd $3$-edge-coloring of $G$, which would give the desired contradiction.)

Begin by noting that the following assertion is equivalent to the existence of $\varphi_i$: \textit{The lobe $G_i$ can be turned into an odd $3$-edge-colorable graph by splitting out $v$ entirely into $2$-vertices and then suppressing all of them except one.} Indeed, if such a splitting modification of $G_i$ exists, then the obtained graph admits an odd edge-coloring with color set $\{1,2,3\}$ so that the colors $2$ and $3$ appear at the $2$-vertex which `survives the suppression', and this coloring readily leads to the desired $\varphi_i$; the other way around, with an existing $\varphi_i$, create the $2$-vertex which remains unsuppressed so that its incident edges are colored by $2$ and $3$ under $\varphi$, and then create the rest of the $2$-vertices so that each of them is incident with equally colored edges.

Therefore, in view of Proposition~\ref{sufficient condition 3-edge-colorability}, it suffices to show that there is an even cycle $C$ in $G_i$ which passes through $v$. Namely, using such a cycle $C$, start the `splitting procedure' by creating a $2$-vertex $v'$ incident to the edges forming $E_C(v)$. Then arbitrarily split out the rest of $v$ into $2$-vertices and suppress them all (except $v'$). The obtained graph is odd $3$-edge-colorable (by Proposition~\ref{sufficient condition 3-edge-colorability}) because $C$ has become an even cycle passing through the $2$-vertex $v'$.

For demonstrating the existence of an even cycle $C\subseteq G$ with $v\in V(C)$, we may assume that every bouquet of $G_i$ incident with $v$ is a singleton. Since $d_{G_i}(v)\geq4$, consider a triplet $x,y,z\in N_{G_i}(v)$. If there is an even $x$-$y$ path in $G_i-v$, we are obviously done. As $G_i-v$ is connected, let $Q$ be an odd $x$-$y$ path. We show there is an even path going from $z$ to the set $\{x,y\}$.  In view of the connectedness, say $P$ is an odd $z$-$x$ path. On a traversal of $P$ from $z$ to $x$, take $w$ to be the first vertex that belongs to $Q$.  Then $zPw\cup wQx$ and $zPw\cup wQy$ are paths (by the choice of $w$). Because $wQx$ and $wQy$ are of opposite parities, we have found an even $z$-$\{x,y\}$ path, which in turn yields an even cycle $C$ through $v$. In conclusion, it is always possible to construct the desired $\varphi_i$.
\hfill$\diamond$

\bigskip

The final piece of our argument is showing that $v$ is of odd degree in precisely one $v$-lobe.

\bigskip

\noindent \textbf{Claim 4.} \textit{There is a unique $j\in\{1,2,\ldots,k\}$ such that $d_{G_j}(v)$ is odd.}

\medskip

\noindent Since $d_G(v)$ is odd, so is $d_{G_j}(v)$ for some $j$. Suppose there are at least two such indices, say $j_1$ and $j_2$. Let $G'=G_{j_1}\cup G_{j_2}$ and $G''=\bigcup\{G_i:i=1,2,\ldots,k \text { and } i\neq j_1,j_2\}$. There exists an odd edge-coloring $\varphi'$ of $G'$ with color set $\{1,2,3\}$ such that $E_G'(v)$ is colored by $1$ and $2$. Indeed, simply take an odd co-forest $H_k$ of $G_{j_k}-v$ for $k=1,2$, color $H_1\cup H_2$ with $3$, and take an edge-coloring of $\hat{H_k}=G_{j_k}-E(H_k)$ with color set $\{1,2\}$ which is odd away from $v$ and so that the color $k$ is odd at $v$ in $\hat{H_k}$. Combine this $\varphi'$ with an edge-coloring $\varphi''$ of $G''$ with color set $\{1,2,3\}$ that is odd away from $v$ and colors $E_{G''}(v)$ by $1$. Such $\varphi''$ can be obtained as follows: take an odd co-forest in each $G_i-v$ for $i\neq j_1,j_2$, and color its edge set by $3$; then color $E_{G''}(v)$ by $1$ and extend to an edge-coloring with color set $\{1,2\}$ (of the still uncolored part of $G''$) such that the coloring is odd away from $v$. We thus obtain an odd $3$-edge-coloring $\varphi'\cup \varphi''$ of $G$, a contradiction.
\hfill$\diamond$

\bigskip

This completes the verification of $(ii)\Rightarrow(i)$, and thus settles $(P)$ and the proposition.
\end{proof}



We are left with the task of determining which $2$-connected graphs belong to $\mathcal{S}_4$. Let us first narrow down the search by establishing a necessary condition.

\begin{proposition}
    \label{block graph}
If $G\in \mathcal{S}_4$ is a block graph, then $G$ can be obtained from a bipartite $2$-connected odd graph by a single edge subdivision.
\end{proposition}
\begin{proof}
Since $G\in \mathcal{S}_4$, its order is odd (by Proposition~\ref{even order}). Hence, there exists a $2$-vertex $v\in V(G)$.  As $G$ is a block graph, every pair of vertices lie on a cycle. Thus, in view of Proposition~\ref{sufficient condition 3-edge-colorability}, vertex $v$ is the only $2$-vertex in $G$. Moreover, every cycle through $v$ is odd. We are left to show that every cycle that avoids $v$ is even.

Letting $N_G(v)=\{u,w\}$, suppose there is an odd cycle $C_o$ in $G-v$. By the $2$-connectedness of $G$, there exist two disjoint $\{u,w\}$-$V(C_o)$ paths, say a $u$-$u'$ path $P$ and a $w$-$w'$ path $Q$. Let $R$ denote the $u'$-$w'$ path along $C_o$ which is even (resp. odd) if $P$ and $Q$ have same (resp. opposite) parities. Then $P\cup R\cup Q$ is an even $u$-$w$ path in $G-v$. However, we have already established that no cycle through $v$ is even. This contradiction proves our point. Namely, the graph $G\%v$, obtained from $G$ by suppressing $v$, is a $2$-connected bipartite odd graph.
\end{proof}


 Not every graph which can be obtained from a bipartite $2$-connected odd graph by a single edge subdivision belongs in $\mathcal{S}_4$. For example, it can be readily seen from Figure~\ref{fig:k35} that the graph obtained from $K_{3,5}$ by a single edge subdivision is odd $3$-edge-colorable. On the other hand, in view of Theorem~\ref{thm:characterizationsubcubic}, this is not the case for the analogous graph obtained from $K_{3,3}$. So our next step is to characterize which graphs obtainable from a bipartite $2$-connected odd graph by a single edge subdivision are members of $\mathcal{S}_4$.

 \begin{figure}[ht!]
	$$
		\includegraphics[scale=0.8]{K35}
	$$
	\caption{An odd $3$-edge-coloring of the graph obtained from $K_{3,5}$ by a single edge subdivision.}
	\label{fig:k35}
\end{figure}



  Consider a graph $G$ that is obtainable from a bipartite $2$-connected odd graph by a single edge subdivision. Thus $\kappa(G)=\kappa'(G)=\delta(G)=2$. Assume $G$ is not essentially $3$-edge-connected, that is, let there be an (edge) $2$-cut $[X,\bar{X}]$ such that $|X|,|\bar{X}|\geq2$. If the unique $2$-vertex falls in $X$, consider the graphs $G'$ and $G''$ constructed as follows: $G'$ is derived from $G[X]$ by adding an edge between the endpoints of the $2$-cut in $X$; $G''$ is obtained from $G[\bar{X}]$ by introducing a new vertex and joining it with the endpoints of the $2$-cut in $\bar{X}$ (equivalently, $G''=G/X$). Notice that both $G',G''$ are obtainable from bipartite $2$-connected odd graphs by single edge subdivisions (see Figure~\ref{fig:konstrukcija}).
  Reversing the process, let $G',G''$ be disjoint graphs that are obtained from $2$-connected bipartite odd graphs by single edge subdivisions. Break a selected edge of $G'$ into two half-edges, then remove the $2$-vertex of $G''$ along with its incident half-edges, and finally glue the half-edges emanating from $G'$ with the half-edges emanating from $G''$ so that two new (whole) edges are created. Call this process \textit{gluing} of $G'$ and $G''$. Notice that the  resulting graph $G$ is also obtainable from a $2$-connected bipartite odd graph by single edge subdivision, and moreover has a non-trivial $2$-cut.
\begin{figure}[ht!]
	$$
		\includegraphics[scale=0.9]{konstrukcija}
	$$
	\caption{Graph $G$ (left), and graphs $G',G''$ (right). The fat edges form the symmetric difference $E(G)\oplus(E(G')\cup E(G''))$.}
	\label{fig:konstrukcija}
\end{figure}

The importance of transforming $G$ into the pair $G',G''$ and vice versa comes from the following.

\begin{proposition}
    \label{glue}
Let $G$ be a graph that is obtained from a $2$-connected bipartite odd graph by a single edge subdivision, and let $[X,\bar{X}]$ be a nontrivial $2$-cut in $G$ such that the unique $2$-vertex is in $X$. With $G',G''$ as described above, the following equivalence holds:
\begin{equation*}
\chi'_o(G)=4 \quad\text{ if and only if }\quad \chi'_o(G')=\chi'_o(G'')=4\,.
\end{equation*}
\end{proposition}
\begin{proof}
Assuming $\chi'_o(G)=4$, we argue by contradiction that $\chi'_o(G')=\chi'_o(G'')=4$. In view of Proposition~\ref{4-edge-colorability}, more than four colors are never required. Let $e\in E(G')\backslash E(G)$. Suppose $\chi'_o(G')\leq3$ and consider an odd $3$-edge-coloring $\varphi$ of $G'$. Extend the restriction $\varphi_{|E(G[X])}$ to $E(G)$ by using the color $\varphi(e)$ for $E(G)\backslash E(G[X])$. This gives an odd $3$-edge-coloring of $G$, a contradiction.

Suppose now that $\chi'_o(G'')\leq3$. We already know from Proposition~\ref{block graph} (or from Proposition~\ref{even order}) that the graph $G'\%e$, obtained from $G'$ by introducing a $2$-vertex on the edge $e$, is odd $3$-edge-colorable. However then an odd $3$-edge-coloring of $G$ arises by combining an odd $3$-edge-coloring of $G'\%e$ and an odd $3$-edge-coloring of $G''$ with the same color set (after possibly permuting colors in the latter). This contradiction settles that from $\chi'_o(G)=4$ it follows that  $\chi'_o(G')=\chi'_o(G'')=4$.

Let us show the reversed implication by contrapositive. Assume an odd $3$-edge-coloring of $G$ exists and consider its restrictions over the $2$-cut $[X,\bar{X}]$. If the two edges from the cut are colored the same, then we have an odd $3$-edge-coloring of $G'$. Contrarily, if the two edges are colored differently, then an odd $3$-edge-coloring of $G''$ readily appears.
\end{proof}

%Reversing the process, let $G',G''$ be disjoint graphs that are obtained from $2$-connected bipartite odd graphs by single edge subdivisions. Break a selected edge of $G'$ into two half-edges, then remove the $2$-vertex of $G''$ along with its incident half-edges, and finally glue the half-edges emanating from $G'$ with the half-edges emanating from $G''$ so that two new (whole) edges are created. Call this process \textit{gluing} of $G'$ and $G''$. Notice that the  resulting graph $G$ is also obtainable from a $2$-connected bipartite odd graph by single edge subdivision, and moreover has a non-trivial $2$-cut.

\section{Characterization of $\mathcal{S}_4$}

As a result of Proposition~\ref{block tree}, we may confine to $2$-connected graphs. Let us denote by $\mathcal{F}$ the family defined inductively as follows:

\begin{itemize}
\item[$(a)$] every Shannon triangle of type $(2,1,1)$ and minimum degree $2$ belongs to $\mathcal{F}$;
\item[$(b)$] every graph that can be obtained from a $2$-connected bipartite cubic graph by a single edge subdivision belongs to $\mathcal{F}$;
\item[$(c)$] every other graph in $\mathcal{F}$ can be constructed by taking disjoint members $G',G''\in\mathcal{F}$ and gluing them together.
\end{itemize}

\smallskip

Clearly, every member of $\mathcal{F}$ is a $2$-connected graph and a subdivision of an odd graph. Note in passing that parts $(a)$ and $(b)$ of the above constructive definition are not entirely `disjoint' as the smallest Shannon triangle of type $(2,1,1)$ is included in both. The family $\mathcal{F}$ happens to be vital for our desired characterization. Namely, it turns out that a block graph $G$ belongs to $\mathcal{S}_4$ if and only if it belongs to $\mathcal{F}$, which we prove next.

 %First we prove that $\mathcal{F}\subseteq\mathcal{S}_4$. The demonstration that every remaining subdivision of a $2$-connected odd graph is odd $3$-edge-colorable is a bit more complex and shall be included  afterwards, see Theorem~\ref{2-connected}.

%\begin{proposition}
%    \label{F}
%If $G\in\mathcal{F}$ then $\chi'_o(G)=4$.
%\end{proposition}
%\begin{proof}
%In view of the equality~\eqref{eqn:shan} and Proposition~\ref{glue}, we only need the following fact: \textit{Every graph that may be obtained from a bipartite cubic graph by a single edge subdivision is not odd $3$-edge-colorable.} The proof of this is a straightforward double-counting argument (see~\cite{AtaPetSkr16} for the details).
%\end{proof}
%
%The last piece of information that clarifies which block graphs belong to $\mathcal{S}_4$ is provided by the next result.
%
%A straightforward consequence of Propositions~\ref{block graph}--\ref{F} and Theorem~\ref{essentially} is the following.

\begin{theorem}
    \label{2-connected}
Let $G\in \mathcal{S}$ be a 2-connected graph. Then the following statements are equivalent:
\begin{itemize}
\item[$(i)$] $G\in \mathcal{F}$;
\item[$(ii)$] $G\in \mathcal{S}_4$.

\end{itemize}
\end{theorem}

\medskip

To prove Theorem~\ref{2-connected}, we proceed by separately establishing that $(i)\Rightarrow (ii)$ and $(ii)\Rightarrow (i)$. Since $\mathcal{F}\subseteq \mathcal{S}$, the former implication amounts to showing that every graph $G\in\mathcal{F}$ has odd chromatic index $\chi'_o(G)=4$. So, in view of the equality~\eqref{eqn:shan} and Proposition~\ref{glue}, we only need to use the following fact: \textit{Every graph that may be obtained from a bipartite cubic graph by a single edge subdivision is not odd $3$-edge-colorable.} The proof of this is a straightforward double-counting argument (see~\cite{AtaPetSkr16} for the details).

\smallskip

The key ingredient for proving the implication $(ii)\Rightarrow (i)$ is provided by the next auxiliary result, that clarifies which $2$-connected and essentially $3$-edge-connected subdivisions of odd graphs require four colors for an odd edge-coloring.

\medskip

\begin{lemma}
    \label{lema}
If $G\in\mathcal{S}_4$ is a $2$-connected and essentially $3$-edge-connected graph, then its maximum degree $\Delta(G)$ or its order $n(G)$ equals $3$.
\end{lemma}

\medskip

\noindent \textit{Proof of Lemma~4.2.} In view of Proposition~\ref{block graph}, let $v$ be the unique $2$-vertex of $G$, and set $N_G(v)=\{u,w\}$. As $\chi'_o(G)=4$, it follows from Proposition~\ref{even order} that the order $n(G)$ is odd. Moreover, by Proposition~\ref{block graph}, the graph $G\%v$ is bipartite, that is, every cycle through $v$ is odd and every cycle avoiding $v$ is even. Since, apart from $v$, every other vertex in $G$ is of odd degree we have the following.

\bigskip

\noindent \textbf{Claim 1.} \textit{No connected even subgraph $H\subseteq G$ is such that $v\in V(H)$ and $n(H)$ is even.}

\medskip

\noindent Arguing by contradiction, note that in the edge-complement $\hat{H}=G-E(H)$, the vertex $v$ is isolated whereas every other vertex has an odd degree. Take an odd factor $K$ of $H$, and color $E(K)$ by $1$, $E(H)\backslash E(K)$ by $2$ and $E(G)\backslash E(H)$ by $3$. This gives an odd $3$-edge-coloring of $G$, a contradiction.
\hfill$\diamond$

\bigskip

In particular, it follows from Claim~1 that there is no pair of cycles $C,C'$, one of which passes through $v$, such that $V(C)\cap V(C')$ is a singleton; call this formation a \textit{forbidden cycle pair} (cf. Figure~\ref{fig:forbidden}).  Several structural constraints arise from the absence of forbidden cycle pairs.

\begin{figure}[ht!]
	$$
		\includegraphics{forbidden}
	$$
	\caption{A forbidden cycle pair.}
	\label{fig:forbidden}
\end{figure}

\smallskip

Supposing $\Delta(G)\geq5$ and $n(G)\geq5$, consider an arbitrary `large' vertex $z$, that is, a vertex of degree $d_G(z)\geq5$. By the $2$-connectedness of $G$, there exists a cycle $C\subseteq G$ such that $v,z\in V(C)$ and $|V(C)|\geq5$. Indeed, if $z\neq u,w$ then any cycle through $v$ and $z$ works; contrarily, select a vertex from $V(G)\backslash\{u,v,w\}$ and use a cycle passing through $v$ and that vertex. Let $P=C-v$ be the $u$-$w$ path that goes through $z$ and is contained within $C$ (it may happen that $z$ is an endpoint of $P$). We consider the collection $\mathcal{P}_z$ of paths $Q$ in $G-v$ such that $P\cap Q$ consists of two isolated points, one of which is $z$. Let us refer to the other endpoint of $Q\in\mathcal{P}_z$ as its \textit{ending}. We denote by $\mathrm{In}(Q)$ the set of \textit{internal vertices} of the path $Q$, those that are not its endpoints.\footnote{Not to be confused with `internal vertices of a block'.}


\bigskip

\noindent \textbf{Claim 2.} \textit{There is a mapping from $E_G(z)\backslash E_C(z)$ to $\mathcal{P}_z$, sending $e\mapsto Q_e$, such that $e$ is the initial edge of $Q_e$ (on a traversing of $Q_e$ that starts at $z$). Moreover, for every such mapping it holds that}

\begin{equation*}
e\neq e' \quad \Rightarrow \quad Q_e \text{ and } Q_{e'} \text{ are internally disjoint, i.e., } \mathrm{In}(Q_e)\cap \mathrm{In}(Q_{e'})=\emptyset\,.
\end{equation*}

\medskip

\noindent If the other endpoint of $e$ (besides $z$) belongs to $V(P)$, then the path $Q_e$ is uniquely determined: namely, it is the $1$-path with edge set $\{e\}$. Contrarily, if that endpoint of $e$ does not belong to $V(P)$, then it falls into a component, say $K_e$, of $G-V(C)$. Note that then  $e\neq e'$ implies $K_e\neq K_{e'}$, for otherwise a forbidden cycle pair (that includes $C$) is present (cf. Figure~\ref{fig:cycle pair}). Thus, $Q_e$ exists in this case as well. (However, as $e\notin\{vu,vw\}$, the ending of $Q_e$ is no longer uniquely determined: indeed, for then $|\partial(K_e)|\geq3$ and no two edges from the edge cut $\partial(K_e)$ have the same endpoint on $P$.)

\begin{figure}[ht!]
	$$
		\includegraphics[scale=0.75]{cyclepair}
	$$
	\caption{A forbidden cycle pair (at $z$) if $K_e=K_{e'}$. Namely, if $x$ and $y$ are the other endpoints (besides $z$) of $e$ and $e'$, respectively, then an $x$-$y$ path taken within the shared component combines with $e$ and $e'$ to produce a cycle.}
	\label{fig:cycle pair}
\end{figure}

 As for the implication $e\neq e' \Rightarrow \mathrm{In}(Q_e)\cap \mathrm{In}(Q_{e'})=\emptyset$, note that it follows immediately from the above discussion.
\hfill$\diamond$

\bigskip

Any subsequent use of notation $Q_e$ is to be understood in the context of Claim~2. We study next the following situation: $e,e'\in E_G(z)\backslash E_C(z)$ are distinct edges and $Q_e,Q_{e'}$ have endings on the same side of $P$ in respect of $z$.

\bigskip

\noindent \textbf{Claim 3.} \textit{Let $Q',Q''\in \mathcal{P}_z$ be internally disjoint, and with their respective endings lying on the same side of $z$ along $P$. If $\bar{z}\in N_P(z)$ is on that side of $P$, then both $Q',Q''$ are $z$-$\bar{z}$ edges.}

\medskip

\noindent First we show that $Q'$ and $Q''$ share the same ending. Arguing by contradiction, suppose their respective endings, say $z'$ and $z''$, differ. Without loss of generality, let $z'$ be an internal vertex of the subpath $zPz''$. Denote $C'=zPz'\cup Q'$ and $C''=zPz''\cup Q''$. Then $C\oplus C''$ and $C'$ constitute a forbidden cycle pair. The obtained contradiction confirms that $Q',Q''$ have the same ending, say $z^*$.

\begin{figure}[ht!]
	$$
		\includegraphics[scale=0.85]{endings}
	$$
	\caption{A detour from $P$ that yields a forbidden cycle pair at $z$.}
	\label{fig:endings}
\end{figure}

Next we prove that $Q',Q''$ are actually $1$-paths. For argument's sake, suppose $\mathrm{In}(Q')\neq\emptyset$. Then, as $\mathrm{In}(Q')\neq\{v\}$, the essential $3$-edge-connectedness of $G$ guarantees that there is an edge $f\notin E(Q')$ such that: $(1)$ it emanates from the component of  $G-V(C)$ containing $\mathrm{In}(Q')$, and $(2)$ has its other endpoint within $V(P)$. However, such an edge $f$ would  contradict the already established feature of shared path endings (either in regard to the large vertex $z$ or to the large vertex $z^*$). Consequently, both $Q'$ and $Q''$ are $z$-$z^*$ edges.

Finally, let us show that $z^*\in N_P(z)$, that is, $z^*=\bar{z}$ is a neighbor of $z$ along $P$. Once again we argue by contradiction and evoke essential $3$-edge-connectedness. Suppose $\mathrm{In}(zPz^*)\neq\emptyset$, and let $H$ be the component of $G-(V(C)\backslash\mathrm{In}(zPz^*))$ that includes $\mathrm{In}(zPz^*)$. Since $|\partial(H)|\geq3$ there exists an edge $h\in \partial(H)\backslash E(P)$. It cannot be that $h$ has an endpoint within $\{z,z^*\}$. Indeed, for otherwise, $h$ would contradict the feature of shared path endings in regard to its endpoint in $\{z,z^*\}$ (see Figure~\ref{fig:neighbor}). Thus, without loss of generality, assume $h$ meets $P$ on the side of $z$ not including $z^*$. If $x$ is the endpoint of $h$ on $P$, then there is a path $R$ that starts at $x$ along $h$ and goes through $H$ until it reaches $P$ again, say at a vertex $y\in \mathrm{In}(zPz^*)$.

\begin{figure}[ht!]
	$$
		\includegraphics[scale=0.85]{neighbor}
	$$
	\caption{A detour from $P$ along $R$ that yields a forbidden cycle pair.}
	\label{fig:neighbor}
\end{figure}

\noindent Define $\bar{C}=xPy \cup R$ and $\bar{\bar{C}}=Q'\cup Q''$. Then $C\oplus \bar{C}$ and $\bar{\bar{C}}$ form a forbidden cycle pair. The obtained contradiction settles the claim.
\hfill$\diamond$

\bigskip

Since $z$ is a large vertex, we have that $|E_G(z)\backslash E_C(z)|\geq3$. Consequently, on at least one side along $P$ the vertex $z$ is incident with a $3^+$-bouquet $\mathcal{B}_{z\bar{z}}$ (see Figure~\ref{fig:bouquet}); call it a \textit{large bouquet}. Thus, every large vertex lying on $P$ is incident with at least one large bouquet (shared with an adjacent large vertex along $P$). Consequently, for each $e\in E_G(z)\backslash E_C(z)$, all the paths $Q_e$ have endings on the same side of $P$ with regard to $z$.

\begin{figure}[ht!]
	$$
		\includegraphics[scale=0.9]{large_bouquet}
	$$
	\caption{A large bouquet at the vertex $z$ (and at the vertex $\overline{z}$).}
	\label{fig:bouquet}
\end{figure}

Concerning the parity of their size, we distinguish between \textit{even} and \textit{odd} large bouquets. In case a large vertex along $P$ is incident with two large bouquets, then clearly one of them is even whereas the other one is odd.

\bigskip

\noindent \textbf{Claim 4.} \textit{If $z',z''$ are adjacent large vertices along $P$ such that $\mathcal{B}_{z'z''}$ is an even large bouquet of $G$, then at least one of them is incident with an odd large bouquet as well.}

\medskip

\noindent Since both $d_G(z'),d_G(z'')$ are odd and $\mathcal{B}_{z'z''}$ is even, the sets $E_G(z')\backslash(E(P)\cup \mathcal{B}_{z'z''})$ and $E_G(z'')\backslash(E(P)\cup \mathcal{B}_{z'z''})$ are even-sized. Suppose neither of $z',z''$ is incident with an odd large bouquet along $P$. Then, in view of Claims~2 and~3, the sets $E_G(z')\backslash(E(P)\cup \mathcal{B}_{z'z''})$ and $E_G(z'')\backslash(E(P)\cup \mathcal{B}_{z'z''})$ are actually empty. However, this gives $|\partial(\{z',z''\})|=2$, contradicting that $G$ is essentially $3$-edge-connected.
\hfill$\diamond$

\bigskip

In particular, as every large vertex along $P$ is incident with at least one large bouquet, there exists an odd large bouquet, say $\mathcal{B}_{z\overline{z}}$. Noting that $\{z,\overline{z}\}\cap\{u,w\}=\emptyset$ (for otherwise every vertex from $\{z,\overline{z}\}\cap\{u,w\}$ would have an even degree $\geq4$), we may assume that $u,z,\overline{z},w$ appear in this relative order along $P$.  We shall gradually construct a triplet $(H,\mathcal{Z},\overline{\mathcal{Z}})$ whose entries are a subgraph $H\subseteq G$ and two subsets $\mathcal{Z}\subseteq V(uPz), \overline{\mathcal{Z}}\subseteq V(\overline{z}Pw)$ so that $\mathcal{B}_{z\overline{z}}\subseteq E(H)$ and the union $\mathcal{Z}\cup\overline{\mathcal{Z}}$ comprises all $3^+$-vertices of $H$. To do so, we initially plug in values $H=C+\mathcal{B}_{z\overline{z}}$, $\mathcal{Z}=\{z\}$ and $\overline{\mathcal{Z}}=\{\overline{z}\}$, and then suitably augment each entry. Once completed, the key features of the subgraph $H$ will be:
\begin{itemize}
\item each of its $3^+$-vertices has equal degrees in $H$ and $G$;
\item it admits an odd edge-coloring $\varphi_H$ with color set $\{1,2,3\}$ such that the color $3$ is used only on an edge of $\mathcal{B}_{z\overline{z}}$.
\end{itemize}

These features of $H$ would immediately imply that the graph $G$ is odd $3$-edge-colorable: namely, as every non-isolated vertex of $\hat{H}$ has an odd degree in $\hat{H}$ and degree $2$ in $H$, simply extend the above mentioned $\varphi_H$ from $E(H)$ to $E(G)$ by coloring $E(G)\backslash E(H)$ with $3$; as $z,\overline{z}$ are isolated in $\hat{H}$, the described extension is an odd $3$-edge-coloring of $G$.

\smallskip

Let us explain the construction of the triplet $(H,\mathcal{Z},\overline{\mathcal{Z}})$. We augment the initial value $H=C+\mathcal{B}_{z\overline{z}}$ along $P$ on each `side' of $\mathcal{B}_{z\overline{z}}$. Since the augmentation process is analogous on both sides, it will only be described for the side of $z$, noting that this half only gives the final value of $\mathcal{Z}$ whereas the final value of $H$ is reached once the process is repeated on the side of $\overline{z}$ as well.



The set $\mathcal{Z}$ is being perceived as linearly ordered (i.e., an ever increasing stack) and comprised of the vertices from $V(H)\cap V(uPz)$ for which $d_H$ exceeds $d_C$. Initially, the head of this stack is the vertex $z$. At each stage, the value $d_G-d_H$ is scanned for the vertex $x$ at the head of $\mathcal{Z}$. The algorithm terminates when this value becomes zero at the head. Meanwhile $d_H(x)<d_G(x)$, and we act as follows:
\begin{itemize}
\item[$(i)$] the subgraph $H$ is augmented with $\bigcup\{Q_e: e\in E_G(x)\backslash E_H(x)\}$, where one path $Q_e$ is added for each edge $e$);
\item[$(ii)$] the shared ending of those $Q_e$'s is appended to $\mathcal{Z}$ as the new head of the stack.
\end{itemize}
\noindent To slightly clarify this two-fold action, with its part $(i)$ the difference $d_G-d_H$ becomes equal to $0$ at the old head, but with its part $(ii)$ the focus has already switched to the new head. We observe in passing that the important fact of the mentioned $Q_e$'s sharing the same ending is guaranteed by Claims~$2$ and $3$. Moreover, depending on the value of $d_G(x)-d_H(x)$, there are two possibilities for $\{Q_e: e\in E_G(x)\backslash E_H(x)\}$, cf. steps 5--11 in the algorithm below. Also note that steps 6 and 7 of this algorithm are justified by Claim~$3$. Namely, in step 6 it is assumed that $d_G(x)-d_H(x)>1$, hence Claim~$3$ assures there is a large bouquet towards $u$ incident with $x$ and its neighbor (see Figure~\ref{fig:steps}). As for steps 9 and 10, there it is assumed that $d_G(x)-d_H(x)=1$, implying $E_G(z)\backslash E_H(z)$ is a singleton. The outcome $(H,\mathcal{Z})$ is not uniquely determined because there is a possible freedom of choice for $Q_e\in\mathcal{P}_x$ in step 9.

\bigskip

\noindent \textbf{Augmentation algorithm} (on the side of $z$)

\medskip

\noindent INPUT: a cycle $C$ through $v$ and an odd large bouquet $\mathcal{B}_{z\overline{z}}$

\smallskip

\noindent OUTPUT: a subgraph $H\subseteq G$ and a queue $\mathcal{Z}\subseteq V(uPz)$

\medskip

\begin{enumerate}

\item[1.] set $t=0$, $\mathcal{Z}=\{z\}$ and $H=C+\mathcal{B}_{z\overline{z}}$\\[-8mm]
\item[2.] consider the head $x$ of $\mathcal{Z}$\\[-8mm]
\item[3.] \textbf{while} $d_H(x)<d_G(x)$ \textbf{do}\\[-8mm]
\item[4.] \quad\quad increment $t$ by $1$\\[-8mm]
\item[5.] \quad\quad \textbf{if} $d_G(x)-d_H(x)>1$ \textbf{then}\\[-8mm]
\item[6.] \quad\quad\quad\quad augment $H$ with the large bouquet $\mathcal{B}$ incident with $x$\\[-8mm]
\item[7.] \quad\quad\quad\quad append the other endpoint of $\mathcal{B}$ to the queue $\mathcal{Z}$\\[-8mm]
\item[8.] \quad\quad \textbf{else}\\[-8mm]
\item[9.] \quad\quad\quad\quad augment $H$ with a path $Q_e\in\mathcal{P}_x$, where $e$ is the unique edge in $E_G(z)\backslash E_H(z)$\\[-8mm]
\item[10.] \quad\quad\quad\quad append the ending of the path $Q_e$ to the queue $\mathcal{Z}$\\[-8mm]
\item[11.] \quad\quad \textbf{end if}\\[-8mm]
\item[12.] \textbf{end while}\\[-8mm]
\item[13.] return $H$ and $\mathcal{Z}$\\[-8mm]

\end{enumerate}

\begin{figure}[ht!]
	$$
		\includegraphics[scale=0.75]{steps}
	$$
	\caption{The augmentation of $H$ and $\mathcal{Z}$ described in steps 6 and 7. The edges currently forming $E(H)\cap E(P)$ are depicted heavier, and the vertices currently forming the stack $\mathcal{Z}$ are depicted as squares. The several light edges in the upper drawing complete the existing large bouquet (towards $u$) incident with the old head $x$ (and also with the new one).}
	\label{fig:steps}
\end{figure}

\bigskip

For the second run of the algorithm, replace every mention of $\mathcal{Z}$ with $\overline{\mathcal{Z}}$, initially defined as $\overline{\mathcal{Z}}=\{\overline{z}\}$. Moreover, on this occasion use as a starting value of $H$ the output value from the first run.
Observe that, as promised, the eventually obtained subgraph $H\subseteq G$ is such that every vertex from $\mathcal{Z}\cup \overline{\mathcal{Z}}$ is isolated in the edge-complement $\hat{H}$, and every other vertex of $H$ has degree $2$. Moreover, the set $E(H)$ has the following important property: it contains at least one odd large bouquet (surely $\mathcal{B}_{z\overline{z}}$ is such), and if all large bouquets are to be removed from $H$, then the remaining cycles would be even and disjoint. In addition, each of them would intersect with either $\mathcal{Z}$ or $\overline{\mathcal{Z}}$ and it would do so in precisely two consecutive vertices (of the corresponding stack), see Figure~\ref{fig:h}.

\begin{figure}[ht!]
	$$
		\includegraphics[scale=0.9]{H}
	$$
	\caption{A sketch of the final outcome for $(H,\mathcal{Z},\overline{\mathcal{Z}})$. All vertices from $\mathcal{Z}\cup \overline{\mathcal{Z}}$ are depicted as squares, and $|\mathcal{Z}|=|\overline{\mathcal{Z}}|=4$. There are five large bouquets within $H$, three of which are odd. If all large bouquets are removed, only two disjoint cycles would remain: one of length $8$  and one of length $10$; the former shares couple of vertices with $\mathcal{Z}$, and the latter with $\overline{\mathcal{Z}}$.}
	\label{fig:h}
\end{figure}

The above features of $E(H)$ enable us to construct the desired $\varphi_H$. We start by coloring $E(C)$ with color set $\{1,2,3\}$ as follows: use $1$ for $uv$, $2$ for $vw$, and $3$ for $z\overline{z}$ (this is the edge in $\mathcal{B}_{z\overline{z}}\cap E(P)$); then by simultaneously traversing $P$ from both endpoints, alternate between the colors $1$ and $2$ at each vertex outside $\mathcal{Z}\cup \overline{\mathcal{Z}}$, and repeat the same color ($1$ or $2$) at each vertex from $(\mathcal{Z}\cup \overline{\mathcal{Z}})\backslash\{z,\overline{z}\}$. Once the described coloring of $E(C)$ is completed, next we color each large bouquet of $H-\mathcal{B}_{z\overline{z}}$ monochromatically, i.e. use for the whole of it the color already appearing on its edge within $P$. As for $\mathcal{B}_{z\overline{z}}$, use the color $1$ on an odd number of its remaining uncolored edges, and use the color $2$ for the rest of it. At this point, the uncolored part of $E(H)$ consists of those $Q_e$'s included in $H$ which are of length $\ell(Q_e)\geq2$. To each of them apply a proper edge-coloring with color set $\{1,2\}$ so that at every endpoint the coloring of $E(H)$ remains locally monochromatic. Such coloring is clearly possible since every such $Q_e$ forms an even cycle with a suitable portion of $P$, and the endpoints of $Q_e$ are consecutive vertices of either $\mathcal{Z}$ or $\overline{\mathcal{Z}}$. \qed

\medskip

The freshly proved Lemma~\ref{lema}, combined with Propositions~\ref{block graph} and~\ref{glue}, settles the implication $(ii)\Rightarrow (i)$. Indeed, arguing by contradiction, let $G$ be a $2$-connected member of $\mathcal{S}_4\backslash \mathcal{F}$ that minimizes the order $n(G)$. Taking into account Propositions~\ref{block graph} and~\ref{glue}, it is implied by part $(c)$ of the construction of $\mathcal{F}$ that $G$ is essentially $3$-edge-connected. Consequently, by Lemma~\ref{lema}, $\Delta(G)=3$ or $n(G)=3$. However, if the former holds then Proposition~\ref{block graph} assures that $G\in\mathcal{F}$ (due to part $(b)$ of the construction). Similarly, if $n(G)=3$ then its assumed $2$-connectedness and the fact that $G\in \mathcal{S}$ imply it is a Shannon triangle of type $(2,2,1)$ with $\delta(G)=2$. Hence $G\in \mathcal{F}$ (due to part $(a)$ of the construction). The obtained contradiction settles that $(ii)\Rightarrow (i)$. This completes the proof of Theorem~\ref{2-connected}.



%--------------------------------------------------------------------------------------------------------------------

\bigskip

  Finally, we arrive at our main result. To formulate it succinctly, the standard notion of a `block-cutpoint graph' proves useful. Recall that for a given connected graph $G$, its \textit{block-cutpoint graph} is a simple bipartite graph in which one partite set consists of the cut-vertices of $G$ and the other has a vertex $b_i$ for each block $B_i$ of $G$; we include $vb_i$ as an edge of the block-cutpoint graph if and only if $v\in B_i$. Our findings can be summarized as follows.



\begin{theorem}
    \label{thm}
Let $G\in \mathcal{S}$ be a connected graph. If $(\mathcal{B},\mathcal{V})$ is the bipartition of the block-cutpoint graph of $G$, where $\mathcal{B}$ is the set of blocks and $\mathcal{V}$ the set of cut-vertices, then the following holds:
\begin{equation*}
\chi'_o(G)=
\begin{cases}
1 & \text{\quad if \,} G \text{ is odd} \,;\\
2 & \text{\quad if \,} G \text{ has 2-vertices, with an even number of them on each cycle}\,;\\
4 & \text{\quad if \,} \mathcal{B}\subseteq\mathcal{F} \text{ and for every }  v\in\mathcal{V} \text{ there is a unique } B\in\mathcal{B} \text{ with odd } d_B(v)\,;\\
3 & \text{\quad otherwise}\,.
\end{cases}
\end{equation*}
\end{theorem}
\begin{proof}
Straightforward from Corollary~\ref{odd 2-edge-colorability}, Proposition~\ref{block tree} and Theorem~\ref{2-connected}.
\end{proof}

\section{Further work}

It is implied by Theorem~\ref{thm} that the problem of determining the odd chromatic index of a subdivision of an odd graph is solvable in polynomial time. In view of Theorems~\ref{odd 4-edge-colorability} and~\ref{Kano}, the complexity questions concerning the value of the odd chromatic index of general graphs amount to deciding on their odd $3$-edge-colorability. As already mentioned in Section~1.2, this question is still open.

We propose the study of a related set of questions arising from the following line of reasoning. The class $\mathcal{S}$ can be captured by using the notion of \textit{maximum even degree}, defined as follows. Let $\Delta_{\mathrm{even}}(G)$ denote the maximum even value among the vertex degrees of $G$. Thus $\mathcal{S}=\{G: G \text{ is a graph and }\Delta_{\mathrm{even}}(G)\leq2\}$, and Theorem~\ref{thm} tells that the problem of determining $\chi'_o(G)$ whenever $\Delta_{\mathrm{even}}(G)\leq2$ can be efficiently solved. For every $k=1,2,\ldots,$ introduce notation $\mathcal{S}^{(k)}=\{G: G \text{ is a graph and }\Delta_{\mathrm{even}}(G)\leq2k\}$. So, by ignoring isolated vertices, $\mathcal{S}^{(0)}=\mathcal{O}$; and obviously  $\mathcal{S}^{(1)}=\mathcal{S}$. We find the next question interesting.

\begin{question}
    \label{q1}
Is the decision problem whether a graph $G\in \mathcal{S}^{(2)}$ has $\chi'_o(G)\leq3$ solvable in polynomial time?
\end{question}

A positive answer for Question~\ref{q1} would open
the door to consider the following more general problem.

\begin{question}
    \label{q2}
Given a positive integer $k$, is the decision problem whether a graph $G\in \mathcal{S}^{(k)}$ has $\chi'_o(G)\leq3$ solvable in polynomial time?
\end{question}

Another possible field of study is to consider $\mathcal{S}^{(k)}$-edge-colorability of graphs for a fixed positive integer $k$ (instead of odd edge-colorability); that is, define a new type of edge-coloring by requiring that each color class is a member of $\mathcal{S}^{(k)}$ (rather than of $\mathcal{O}=\mathcal{S}^{(0)}$). Say the corresponding index (representing the minimum sufficient number of colors) is $\chi'_{\mathcal{S}^{(k)}}(G)$. For example, it is readily observed that $\chi'_{\mathcal{S}}(W_4)=2$ as opposed to $\chi'_o(W_4)=4$. Similarly, if $G$ is a Shannon triangle of type $(2,2,2)$, then $\chi'_{\mathcal{S}}(G)\leq3$ in contrast to $\chi'_o(G)=6$. Note that there are graphs that require at least four colors for an $\mathcal{S}$-edge-coloring. Namely, every Shannon triangle $G$ of type $(2,2,1)$ and multiplicity at least $3$ has $\chi'_{\mathcal{S}}(G)=4$ (see Figure~\ref{fig:s}).

\begin{figure}[ht!]
	$$
		\includegraphics[scale=1.1]{shanno}
	$$
	\caption{A Shannon triangle of type $(2,2,1)$ that requires four colors for an $\mathcal{S}$-edge-coloring. Its bouquets are of size $4, 4$ and $3$, respectively.}
	\label{fig:s}
\end{figure}

 We are tempted to end our discussion here with the following.

%\begin{problem}
%    \label{q3}
%Given a positive integer $k$, characterize connected loopless graphs $G$ that have $\chi'_{\mathcal{S}^{(k)}}(G)\leq3$.
%\end{problem}
%
%In particular, we believe the following holds true.

\begin{conjecture}
    \label{conj:3}
If $G$ is a connected loopless graph that is not a Shannon triangle of type $(2,2,1)$ and multiplicity $\geq3$, then $\chi'_{\mathcal{S}}(G)\leq3$.
\end{conjecture}



%In view of Theorem~\ref{Kano}, our discussion here ends with another complexity question.
%
%\begin{question}
%    \label{q4}
%Is the decision problem whether a given graph $G$ has $\chi'_{\mathcal{S}}(G)\leq2$ solvable in polynomial time?
%\end{question}

%----------------------------------------------------------------------------------------------------------------------
%----------------------------------------------------------------------------------------------------------------------


\bigskip
\noindent
{\bf Acknowledgements.}
This work is partially supported by ARRS Program P1-0383 and ARRS Project J1-1692.


\bibliographystyle{plain}

\begin{thebibliography}{99}

\bibitem{AtaPetSkr16} R.~Atanasov, M.~Petru\v{s}evski, R.~\v{S}krekovski, \textit{Odd edge-colorability of subcubic graphs}, Ars Math. Contemp. $\mathbf{10}$ (2016), 359--370.
\bibitem{BonMur08} J.~A.~Bondy, U.~S.~R.~Murty, \textit{Graph Theory}, Graduate Texts in Mathematics, Springer, New York $\mathbf{244}$ (2008).
%\bibitem{KanKat02} M.~Kano, G.~Y.~Katona, \textit{Odd subgraphs and matchings}, Discrete Math. $\mathbf{250}$ (2002), 265--272.
\bibitem{KanKatVar18} M.~Kano, G.~Y.~Katona, K.~Varga, \textit{Decomposition of a graph into two disjoint odd subgraphs}, Graphs and Combin. $\mathbf{34}(6)$ (2018) 1581--1588.
\bibitem{LuzPetSkr15} B.~Lu\v{z}ar, M.~Petru\v{s}evski, R.~\v{S}krekovski, \textit{Odd edge coloring of graphs}, Ars Math. Contemp. $\mathbf{9}$ (2015)  277--287.
\bibitem{Mat06} T.~Matrai, \textit{Covering the edges of a graph by three odd subgraphs}, J. Graph Theory $\mathbf{53}$ (2006) 75--82.
\bibitem{Pet18} M.~Petru\v{s}evski, \textit{Odd $4$-edge-colorability of graphs}, J. Graph Theory $\mathbf{87}(4)$ (2018) 460--474.
\bibitem{PetSkr21} M.~Petru\v{s}evski, R.~\v{S}krekovski, \textit{Odd decompositions and coverings of graphs}, European J. Combin. $\mathbf{91}$ (2021) 103225.
\bibitem{Pyb91} L.~Pyber, \textit{Covering the edges of a graph by...}, Sets, Graphs and Numbers, Colloquia Mathematica Societatis J\'{a}nos Bolyai $\mathbf{60}$ (1991), 583--610.
\bibitem{Sch03} A.~Schrijver, \textbf{Combinatorial optimization. {P}olyhedra and efficiency.} {V}ol. {A}, Algorithms and Combinatorics, Springer-Verlag, Berlin (2003).



\end{thebibliography}


\end{document} 